\chapter{Your first assembly program}
\MARSChapterHeading{Your first assembly program}
\label{revision:first-program}

A program that reads a number, doubles it and adds one is easy to describe in a
language you already know. Assembly makes the intermediate work explicit. You
choose where the input is held, which instruction computes each result, and how
the result reaches the output operation. We will follow that small calculation
from source code to an observable change in the machine.

The program accepts an integer between $-1000$ and $1000$. For input 20, it prints
41. For input $-3$, it prints $-5$. This range keeps the arithmetic comfortably
within a signed 32-bit word. Wider values and overflow will receive a fuller
treatment after we have a working model of execution.

By the end of this chapter, you should be able to explain where a value comes
from, which instruction changes it, and where it goes next. You will also change
the calculation and use a debugger to investigate a wrong result.

\section{The state that instructions change}

A processor has a small collection of storage locations called \emph{registers}.
An instruction can read values from registers, perform an operation and write a
result to a register. Register names in this assembly language begin with a dollar
sign: \texttt{\$v0} and \texttt{\$t0} name two different locations. For this first
program, you need only a few names; there is no reason to memorize the complete
register set before using it.

Consider input 20. An input operation places 20 in \texttt{\$v0}. An addition reads
that register twice, adds the values and puts 40 in \texttt{\$t0}. A second addition
replaces 40 in \texttt{\$t0} with 41. The original input and the intermediate result
occupy different registers. A source register is read because the instruction names
it, not because it happens to contain a recent value.

Registers are distinct from the larger byte-addressed \emph{memory} that holds
program instructions and data. Our prompt is stored as bytes in memory. The input
number is returned in a register. We will use both kinds of storage here, then
study their relationship in detail in the memory chapter.

\section{Reading the source}

The complete example is shown below. Read it once to identify the three parts:
data declarations, the input and calculation, and output followed by termination.
The next sections explain those parts. You do not need to infer the meaning of
all the lines before proceeding.

\MARSCodeLabel{A complete console program}
\begin{minted}{mars-idea}
.data
prompt:  .asciiz "Enter a small integer: "
answer:  .asciiz "Twice your number plus one is: "
newline: .asciiz "\n"

.text
.globl main
main:
    la $a0, prompt
    li $v0, 4
    syscall
    li $v0, 5
    syscall
calculate:
    addu $t0, $v0, $v0
    addiu $t0, $t0, 1
report:
    la $a0, answer
    li $v0, 4
    syscall
    move $a0, $t0
    li $v0, 1
    syscall
    la $a0, newline
    li $v0, 4
    syscall
done:
    li $v0, 10
    syscall
\end{minted}

An \emph{assembler} translates this source into instructions and data that MARS
can execute. Lines beginning with a dot give instructions to the assembler.
\texttt{.data} introduces data declarations and \texttt{.text} introduces executable
code. \texttt{.globl main} makes the entry name available outside its local definition.

A name followed by a colon is a \emph{label}. It names a location; it does not run
an operation. Here, \texttt{main} names the beginning of the program, while
\texttt{calculate} and \texttt{report} make useful places to find in the debugger.
The three labels in the data section name the beginnings of strings.
\texttt{.asciiz} stores the string's bytes followed by a zero byte marking its end.
The sequence \texttt{\textbackslash n} represents a newline in a string.

The remaining lines contain instructions. The first word is the operation name,
called its \emph{mnemonic}; the names and constants after it are its operands.
Some convenient forms, including \texttt{li}, \texttt{la} and \texttt{move}, are
assembler pseudoinstructions. The assembler may translate one source line into
more than one machine instruction. For now, follow the source-level calculation;
the debugger can also expose the underlying instructions when you need them.

\section{Asking MARS for input and output}

Arithmetic instructions operate on values already available to the processor.
Console input and output require another interface. MARS supplies an educational
service convention: put a service number in \texttt{\$v0}, place any arguments in
the designated registers, and execute \texttt{syscall}. These service numbers are
part of the MARS environment, not a universal convention for every MIPS system.

The first three instructions ask for string output. \texttt{la \$a0, prompt}
places the prompt's address in \texttt{\$a0}. An address identifies where the string
begins; it is not the string itself. \texttt{li \$v0, 4} selects the print-string
service, and \texttt{syscall} performs the request. The console now displays the
prompt, and execution proceeds to the next instruction.

Next, \texttt{li \$v0, 5} selects integer input. This replaces the previous service
number in that register. The following \texttt{syscall} waits for an integer and
returns it in \texttt{\$v0}. The register has changed roles: it held a request code
before the call and holds the input afterward. For input 20, that returned value
is the number the arithmetic will read.

An interface is an agreement about such roles. The processor does not decide that
20 must be the input because it looks like a plausible number. The service and
our program agree on the register used for the result.

\section{Following the calculation}

The line \texttt{addu \$t0, \$v0, \$v0} names the destination first, then two sources.
With 20 in \texttt{\$v0}, it adds $20+20$ and writes 40 to \texttt{\$t0}. Reading the
same source twice is valid. The operation leaves \texttt{\$v0} unchanged.

The next line, \texttt{addiu \$t0, \$t0, 1}, has a constant operand. Such a constant
is called an \emph{immediate}. The instruction reads the old value of
\texttt{\$t0}, adds one and writes the new value to the same register. Naming one
register as both source and destination does not cause a conflict: the old value
supplies the input and the result replaces it.

These particular additions operate on 32-bit words. On the MIPS64 target, their
word result is sign-extended into the 64-bit destination register. The letter
\texttt{u} indicates that signed overflow does not raise an arithmetic exception;
it does not force negative inputs to become positive. Within our stated range,
the results agree with ordinary integer addition.

\begin{center}
\begin{tabular}{lll}
\toprule
Execution boundary & \texttt{\$v0} & \texttt{\$t0} \\
\midrule
After integer input & 20 & Not yet relevant \\
After doubling & 20 & 40 \\
After adding one & 20 & 41 \\
\bottomrule
\end{tabular}
\end{center}

For input $-3$, the same operations produce $-6$ and then $-5$. Try that trace on
paper before running it. A prediction gives you a specific claim to compare with
the machine, rather than a screen full of values to inspect without a question.

\section{Keeping a result until it is used}

The calculation is complete, but output still requires several requests. The
program first prints the answer prefix, then prints the number, then a newline.
Each request needs another service number in \texttt{\$v0}. The answer therefore
remains in \texttt{\$t0} while the output registers are prepared.

\texttt{move \$a0, \$t0} copies the number into the argument register for integer
output. It does not erase the source. Loading service number one into
\texttt{\$v0} and executing \texttt{syscall} prints that number. Selecting service
four again lets the program print the newline string. Finally, service ten ends
the program normally.

This is the first example of a value's \emph{lifetime}: the answer must survive
from the instruction that produces it until the operation that consumes it.
Keeping the answer in a register that is immediately reused for another purpose
would lose it. Later, function calls will make these lifetime questions more
important, and we will give them explicit conventions.

\section{Using the debugger to answer a question}

Run the supplied example and enter 20. You should see 41 after the answer prefix.
That checks the final result for one input. Debugging lets you investigate the
intermediate reasoning as well.

Set a breakpoint at the addition under \texttt{calculate} and start Debug. Enter
20 when the console asks. When execution pauses before that addition, the returned
input should already be in \texttt{\$v0}. The highlighted instruction is the next
one to execute. Predict its destination value, step once, and inspect
\texttt{\$t0}. Repeat for the following addition. Check whether the debugger is
showing signed decimal or hexadecimal before interpreting a negative result.

If the printed answer is wrong, start with the last boundary where the value was
correct. A correct input followed by a wrong doubled value points to the first
addition. A correct 41 in \texttt{\$t0} followed by wrong output points to the
argument preparation or service selection. This divides the investigation into
questions you can actually test.

\section{Writing a related program}

Change the arithmetic to calculate $3x+1$. Keep the input and output interface the
same. Work out the intermediate values for 20, zero and $-3$, then run all three
cases. You should be able to explain which instructions consume the original
input and which consume an intermediate result.

As a separate experiment, change the integer-output argument preparation to copy
\texttt{\$v0} instead of \texttt{\$t0}. Predict what the console will print by
finding the most recent instruction that wrote \texttt{\$v0}; do not assume that
it still holds the original input. Use a breakpoint to show that the arithmetic
can be correct while the output argument is wrong, then repair the program.

The companion lab develops these observations interactively. Its written
explanations and the program remain available so you can revisit the reasoning
without replaying every IDE action.
